\section*{第 \S\arabic{exercisegroup} 节习题}
\phantomsection
\addcontentsline{toc}{section}{第 \protect\S\arabic{exercisegroup} 节习题}

\begin{enumerate}
\def\labelenumi{\arabic{enumi})}
\item
  设 E 是 R 上的一个完备赋范空间，F 是一个拓扑空间，J 是 R 的一个不退化为一点的区间；设 \((t, \xi) \mapsto A(t, \xi)\) 是一个由 \(J \times F\) 到 \(\mathcal{L}(E)\) 中的映射，使得当 \(\xi\) 趋于 \(\xi_{0}\) 时，\(A(t, \xi)\) 在 J 上一致趋于 \(A(t, \xi_{0})\) 。若 \(C(t, t_{0}, \xi)\) 是线性方程 \(d\mathbf{x}/dt = A(t, \xi)\) . x 的预解，证明：对每个紧区间 \(K \subset J\) ，\(C(t, t_{0}, \xi)\) 一致趋于 \(C(t, t_{0}, \xi_{0})\) 于 \(K \times K\) 上（当 \(\xi\) 趋于 \(\xi_{0}\) 时）。
\item
  设 \(t \mapsto A(t)\) 是一个由 \(J\) 到 \(\mathcal{L}(E)\) 中的规则映射，使得对任意两点 \(s, t\) （在 \(J\) 中），\(A(s)\) 与 \(A(t)\) 可交换。令 \(B(t) = \int_{t_0}^{t} A(s) ds\) 。证明预解 \(C(t, t_0)\) ——方程 \(d\mathbf{x}/dt = A(t)\mathbf{x}\) 的预解——等于 \(\exp(B(t))\) 。
\end{enumerate}

若 \(t \mapsto A_1(t)\) 、\(t \mapsto A_2(t)\) 是两个由 J 到 \(\mathcal{L}(\mathrm{E})\) 中的规则函数，使得对任意两点 \(s\) 、\(t\) （在 J 中），\(A_1(s)\) 、\(A_1(t)\) 、\(A_2(s)\) 、\(A_2(t)\) 两两可交换，证明

\[
\exp \left(\int_ {t _ {0}} ^ {t} \left(A _ {1} (s) + A _ {2} (s)\right) d s\right) = \exp \left(\int_ {t _ {0}} ^ {t} A _ {1} (s) d s\right) \exp \left(\int_ {t _ {0}} ^ {t} A _ {2} (s) d s\right).
\]

\begin{enumerate}
\def\labelenumi{\arabic{enumi})}
\setcounter{enumi}{2}
\tightlist
\item
  给定两个矩阵
\end{enumerate}

\[
A = \left( \begin{array}{c c} 0 & 1 \\ 0 & 0 \end{array} \right), \qquad \qquad B = \left( \begin{array}{c c} 0 & 0 \\ 1 & 0 \end{array} \right)
\]

证明 \(\exp(A+B)\neq\exp(A)\exp(B)\) 。

\begin{enumerate}
\def\labelenumi{\arabic{enumi})}
\setcounter{enumi}{3}
\item
  若 \(P\) 是 \(\mathcal{L}(\mathrm{E})\) 中任意一个自同构，证明 \(\exp(PAP^{-1}) = P\exp(A)P^{-1}\) 对每个 \(A \in \mathcal{L}(\mathrm{E})\) 成立。
\item
  用例子证明：线性方程 \(d\mathbf{x} / dt = A.\mathbf{x} + e^{\alpha t}\mathbf{p}(t)\) （其中 \(A \in \mathcal{L}(\mathrm{E})\) 与 \(t\) 无关，\(\mathbf{p}\) 是一个多项式（系数在 E 中））可以有一个积分等于一个与 \(\mathbf{p}\) 同次数的多项式，即使当 \(\alpha\) 是 \(A\) 的特征根时也是如此。
\item
  设 \(f(X)\) 是一个 n 次复系数多项式，并设
\end{enumerate}

\[
\frac {1}{f (\mathrm{X})} = \sum_ {j = 1} ^ {q} \sum_ {h = 1} ^ {n _ {j}} \frac {\lambda_ {j h}}{(\mathrm{X} - r _ {j}) ^ {h}}
\]

是有理分式 \(1 / f(\mathrm{X})\) 的部分分式分解 (Alg., VII. 7, n° 2)。证明：对任意规则函数 \(b\) ，函数

\[
t \mapsto \sum_ {j = 1} ^ {q} \sum_ {h = 1} ^ {n _ {j}} \lambda_ {j h} \int_ {t _ {0}} ^ {t} \frac {(t - s) ^ {h - 1}}{(h - 1) !} e ^ {r _ {j} (t - s)} b (s) d s
\]

是方程 \(f(\mathrm{D})x = b(t)\) 的一个积分。

\begin{enumerate}
\def\labelenumi{\arabic{enumi})}
\setcounter{enumi}{6}
\tightlist
\item
  \begin{enumerate}
  \def\labelenumii{\alph{enumii})}
  \tightlist
  \item
    设 \(t \mapsto A(t)\) 是一个由区间 \(J \subset R\) 到 \(\mathcal{L}(E)\) 中的映射，使得对每个 \(x \in E\) ，映射 \(t \mapsto A(t).x\) （由 J 到 E 中）连续。证明 \(t \mapsto \|A(t)\|\) 在每个紧区间 \(K \subset J\) 上有界（用 Gen.~Top., IX, p.~194, th. 2）。在这些条件下证明：对每一点 \((t_{0}, \mathbf{x}_{0}) \in \mathbf{J} \times \mathbf{E}\) ，方程 \(d\mathbf{x}/dt = A(t).\mathbf{x} + \mathbf{b}(t)\) 有一个且仅有一个解，它定义在 J 上且等于 \(x_{0}\) 于点 \(t_{0}\) 。进一步，若以 \(\mathbf{u}(t, t_{0}, \mathbf{x}_{0})\) 表示这个解，则映射 \(\mathbf{x}_{0} \mapsto \mathbf{u}(t, t_{0}, \mathbf{x}_{0})\) 是 E 到其自身上的一个双射双连续线性映射 \(C(t, t_{0})\) ，它满足关系式 (6) (IV, p.~179) 与 (7) (IV, p.~181)；进一步，映射 \((s, t) \mapsto C(s, t)\) 是由 \(J \times J\) 到 \(\mathcal{L}(E)\) 中的连续映射。
  \end{enumerate}
\end{enumerate}

\begin{enumerate}
\def\labelenumi{\alph{enumi})}
\setcounter{enumi}{1}
\tightlist
\item
  取 E 为实数序列 \(\mathbf{x}=(x_{n})_{n\in\mathbf{N}}\) （满足 \(\lim_{n\to\infty}x_{n}=0\) ）所成的空间，赋以范数 \(\|x\|=\sup_{n}|x_{n}|\) ，E 在其下是完备的。对每个 \(t\in J=[0,1]\) ，设 \(A(t)\) 是 E 到其自身的、使得
\end{enumerate}

\[
A (t). \mathbf {x} = \left(\frac {1}{1 + n t} x _ {n}\right) _ {n \in \mathbf {N}}.
\]

证明 \(A(t)\) 满足 a) 的条件，但映射 \(t \mapsto A(t)\) （由 J 到 \(\mathcal{L}(\mathrm{E})\) 中）在点 t = 0 处不连续，并且预解 \(C(t, t_0)\) ——方程 \(d\mathbf{x}/dt = A(t)\) . x 的预解——使得映射 \(t \mapsto C(t, t_0)\) （由 J 到 \(\mathcal{L}(\mathrm{E})\) 中）在点 t = 0 处不可导。

* 8) 设 G 是域 R 上的一个有单位元 e 的完备赋范代数。

\begin{enumerate}
\def\labelenumi{\alph{enumi})}
\item
  设 \(t \mapsto \mathbf{a}(t)\) 是区间 \(J \subset R\) 上的一个取值于 G 的规则函数。证明线性方程 \(d\mathbf{x}/dt = \mathbf{a}(t)\mathbf{x}\) 的取值 e 于一点 \(t_0 \in J\) 的积分 u 在 J 上可逆，且其逆是方程 \(d\mathbf{x}/dt = -\mathbf{x}\mathbf{a}(t)\) 的解（若 v 是后一方程的取值 e 于点 \(t_0\) 的积分，考虑 uv 与 vu 所满足的线性方程）。由此推出：对每个 \(x_0 \in G\) ，\(d\mathbf{x}/dt = \mathbf{a}(t)\mathbf{x}\) 的取值 \(x_0\) 于点 \(t_0\) 的积分等于 \(\mathbf{u}(t)\mathbf{x}_0\) 。
\item
  设 \(\mathbf{a}(t)\) 、\(\mathbf{b}(t)\) 是 J 上的两个规则函数，u 与 v 分别是方程 \(d\mathbf{x}/dt = \mathbf{a}(t)\mathbf{x}\) 与 \(d\mathbf{x}/dt = \mathbf{x}\mathbf{b}(t)\) 的取值 e 于点 \(t_{0}\) 的积分。证明方程 \(d\mathbf{x}/dt = \mathbf{a}(t)\mathbf{x} + \mathbf{x}\mathbf{b}(t)\) 的取值 \(x_{0}\) 于点 \(t_{0}\) 的积分等于 \(\mathbf{u}(t)\mathbf{x}_{0}\mathbf{v}(t)\) 。
\item
  设 \(\mathbf{a}(t)\) 、\(\mathbf{b}(t)\) 、\(\mathbf{c}(t)\) 、\(\mathbf{d}(t)\) 是 \(\mathbf{J}\) 上的四个规则函数，\((\mathbf{u}, \mathbf{v})\) 是由两个线性方程组成的方程组
\end{enumerate}

\[
\frac {d \mathbf {x}}{d t} = \mathbf {a} (t) \mathbf {x} + \mathbf {b} (t) \mathbf {y}, \quad \frac {d \mathbf {y}}{d t} = \mathbf {c} (t) \mathbf {x} + \mathbf {d} (t) \mathbf {y}.
\]

证明：若 \(\mathbf{v}\) 在 \(J\) 上可逆，则 \(\mathbf{w} = \mathbf{u}\mathbf{v}^{-1}\) 是方程 \(d\mathbf{z}/dt = \mathbf{b}(t) + \mathbf{a}(t)\mathbf{z} - \mathbf{z}\mathbf{d}(t) - \mathbf{z}\mathbf{c}(t)\mathbf{z}\) （“里卡蒂方程”）的一个积分；反之亦然。由此推出：后一方程在 \(J\) 上的每个取值 \(\mathbf{x}_0\) 于点 \(t_0\) 的积分都具有形式 \((A(t)\mathbf{x}_0 + B(t)\mathbf{e})(C(t)\mathbf{x}_0 + D(t)\mathbf{e})^{-1}\) ，其中 \(A(t), B(t), C(t)\) 与 \(D(t)\) 是由 \(J\) 到 \(\mathcal{L}(E)\) 中的连续映射，满足恒等式 \(U.(xy) = (U.x)y\) 。

\begin{enumerate}
\def\labelenumi{\alph{enumi})}
\setcounter{enumi}{3}
\tightlist
\item
  设 \(w_{1}\) 是 \(d\mathbf{z}/dt = \mathbf{b}(t) + \mathbf{a}(t)\mathbf{z} - \mathbf{z}\mathbf{d}(t) - \mathbf{z}\mathbf{c}(t)\mathbf{z}\) 在 J 上的一个积分；证明：若 w 是这个方程的另一个积分，使得 \(w - w_{1}\) 在 J 上可逆，则 \(w - w_{1}\) 可以用方程 \(d\mathbf{x}/dt = -(\mathbf{d} + \mathbf{c}\mathbf{w}_{1})\mathbf{x}\) 与 \(d\mathbf{x}/dt = \mathbf{x}(\mathbf{a} - \mathbf{w}_{1}\mathbf{c})\) 的积分表出（考虑 \((\mathbf{w} - \mathbf{w}_{1})^{-1}\) 所满足的方程，并用 b)）。
\end{enumerate}

\begin{enumerate}
\def\labelenumi{\arabic{enumi})}
\setcounter{enumi}{8}
\tightlist
\item
  设 \(y_{k}\) ( \(1 \leqslant k \leqslant n\) ) 是 \(n\) 个定义在区间 \(I \subset \mathbf{R}\) 上、具有连续 \((n - 1)^{th}\) 导数的实函数。
\end{enumerate}

\begin{enumerate}
\def\labelenumi{\alph{enumi})}
\item
  证明：若 \(n\) 个函数 \(y_{k}\) 线性相关，则矩阵 \((y_{k}^{(h)}(t))_{0\leqslant h\leqslant n - 1, 1\leqslant k\leqslant n}\) 的秩 \(< n\) （在每一点 \(t\in I\) 处）。
\item
  反之，若对每个 \(t \in I\) 矩阵 \((y_{k}^{(h)}(t))_{0 \leqslant h \leqslant n-1, 1 \leqslant k \leqslant n}\) 的秩 \textless{} n，证明：在每个非空开区间 \(J \subset I\) 中存在一个非空开区间 \(U \subset J\) ，使得 \(y_{k}\) 在 U 上的限制线性相关（若 p 是那些数 q \textless{} n ——使函数 \(y_{k}\) 中任意 q 个的朗斯基行列式在 J 上恒为零的数——中最小的一个，考虑一点 \(a \in J\) ，在该处函数 \(y_{k}\) 中 p - 1 个的朗斯基行列式不为零，并证明函数 \(y_{k}\) 中 p 个在 a 的一个邻域上是一个 p - 1 阶线性方程的积分）。
\item
  令 \(y_{1}(t) = t^{2}\) （对 \(t \in \mathbf{R}\) ），\(y_{2}(t) = t^{2}\) （对 \(t \geqslant 0\) ），\(y_{2}(t) = -t^{2}\) （对 \(t < 0\) ）；证明 \(y_{1}\) 与 \(y_{2}\) 在 \(\mathbf{R}\) 上有连续导数，且 \(y_{1}y_{2}' - y_{2}y_{1}' = 0\) ，但 \(y_{1}\) 与 \(y_{2}\) 在 \(\mathbf{R}\) 上不线性相关。
\end{enumerate}

* 10) 设 \(t \mapsto X(t)\) 是一个由区间 \(\mathbf{J} \subset \mathbf{R}\) 到 \(n\) 阶复矩阵空间中的映射。假设导数 \(t \mapsto X'(t)\) 存在且在 \(\mathbf{J}\) 上连续，并且使得 \(X(t)X'(t) = X'(t)X(t)\) 对每个 \(t \in \mathbf{J}\) 成立。

\begin{enumerate}
\def\labelenumi{\alph{enumi})}
\item
  进一步假设对每个 \(t \in J\) ，\(X(t)\) 的特征值互异。证明这时存在一个常数可逆矩阵 \(P_{0}\) ，使得 \(P_{0}X(t)P_{0}^{-1} = D(t)\) ，其中 \(D(t)\) 是一个对角矩阵，从而 \(X(t_{1})\) 与 \(X(t_{2})\) 对 t 在 J 中的每一对值 \(t_{1}, t_{2}\) 可交换。（在 J 的每一点的一个邻域上写出 \(X(t) = P(t)D(t)P(t)^{-1}\) ，并建立 \(P(t)\) 所满足的微分方程。）
\item
  矩阵
\end{enumerate}

\[
X (t) = \left( \begin{array}{c c c} t ^ {2} & t ^ {3} & t ^ {4} \\ - 2 t & - 2 t ^ {2} & - 2 t ^ {3} \\ 1 & t & t ^ {2} \end{array} \right)
\]

与其导数可交换，但 a) 的结论不成立。*

